% Section 1 — Introduction
% Paper: "The Learnable Bridge: Task Fingerprinting and Adapter Composition
%         via Structured Coupling in Low-Rank Adaptation"

\section{Introduction}
\label{sec:introduction}

A trained LoRA adapter~\cite{hu2022lora} is opaque.
Its two projection matrices $A \in \mathbb{R}^{r \times d}$ and
$B \in \mathbb{R}^{d \times r}$ encode whatever the fine-tuning
objective required, but there is no structured way to inspect
\emph{what} they learned without running inference on held-out data.
Two adapters with identical architectures, trained on different tasks,
are indistinguishable from their weight statistics alone: both are
low-rank, both are small, both produce $\Delta W = BA$.
The adapter is a black box within a black box.

This opacity has practical consequences.
Adapter merging methods---task arithmetic~\cite{ilharco2023editing},
TIES-Merging~\cite{yadav2023ties}, DARE~\cite{yu2024dare},
AdaMerging~\cite{yang2024adamerging}---operate on full weight matrices
or their SVD factors without any compact representation of task identity.
Overfitting detection requires validation loss curves.
Adapter selection from a hub or registry relies on metadata, not on the
adapter's internal state.
In each case, the adapter itself provides no summary of its own content.

We introduce a \emph{bridge matrix}: a learnable $C \times C$ coupling
matrix inserted between $A$ and $B$ in the LoRA factorization.
The down-projection $A$ maps the input to a rank-$r$ hidden vector,
which is reshaped into $C$ channels of width $r/C$; the bridge mixes
channels; the result is reshaped and passed to the up-projection $B$.
With $C = 6$ (motivated by the face-count of the rhombic dodecahedron
in the RhombiLoRA architecture~\cite{rhombilora_paper1, rhombilora_paper2}),
the bridge adds 36 parameters per adapter layer---a 0.03\% overhead
relative to the $\sim$135K parameters in a rank-24 adapter at each
target module.
When the bridge is the identity matrix $I_C$, the architecture
reduces exactly to standard LoRA.

Through systematic analysis of 336 bridge matrices trained across three
tasks (general instruction following, code generation, mathematical
reasoning) on Qwen2.5-7B, we establish three findings:

\paragraph{(a) Task fingerprinting.}
Bridge matrices encode task identity.
A leave-one-out SVM trained on the 28 query-projection bridges
(1,008 parameters total) classifies task type at 84.5\% accuracy,
compared to 33.3\% chance ($p < 10^{-6}$, Mann-Whitney $U$).
Tasks modulate bridge coupling \emph{magnitude}---general instruction
following produces $2.2\times$ more cross-channel coupling than code
generation (Fiedler connectivity 0.040 vs.\ 0.018)---while preserving
coupling \emph{structure}: eigenspectrum cosine similarity exceeds
0.999 between all task pairs.
The fingerprint emerges monotonically from training
(0\% at step 0, 66\% at step 500, 93\% at step 1500 for the hardest
binary pair) with no initialization confound, since all bridges start
as $I_C$.

\paragraph{(b) Smooth adapter composition.}
Linear interpolation of bridge matrices between task-specific adapters
preserves eigenspectrum structure (cosine $> 0.999$) at all mixing
ratios $\alpha \in [0, 1]$, while random-weight bridges retain only
cosine $\approx 0.5$.
The Fiedler trajectory is linear for structurally similar task pairs
(alpaca$\leftrightarrow$code, $R^2 = 0.956$) but non-linear for
dissimilar pairs (code$\leftrightarrow$math, $R^2 = 0.823$), with
destructive interference visible as a connectivity dip at $\alpha
\approx 0.2$---a phenomenon invisible to full-weight merging diagnostics.
This suggests that bridge interpolation can serve as a low-dimensional
probe of adapter compatibility before committing to an expensive merge.

\paragraph{(c) Non-uniform informativeness.}
Query projections carry the most task-discriminative bridge structure
across all three experimental phases: highest coupling (Fiedler 0.050
vs.\ 0.035--0.039 for K/V/O), highest between-task distance, highest
merge sensitivity.
Cross-phase Pearson $r = 0.921$ between task discrimination
(Phase~1A) and merge sensitivity (Phase~2A) confirms that the modules
encoding the most task identity are also the most vulnerable to
cross-task interference during composition.
Later transformer layers are significantly more task-specific than
earlier layers (Spearman $\rho = 0.701$, $p < 0.0001$), with the
output projection surpassing the query projection in the final five
layers (Layer~24 $o_\text{proj}$: distance 0.319, the single most
discriminative position in the model).

\paragraph{The null result.}
The bridge does not improve training loss or perplexity relative to
standard LoRA at matched rank.
The original motivation---that the rhombic dodecahedron's face-pairing
topology would induce direction-specific channel coupling from
structured training data---produced a comprehensive null
(co-planar/cross-planar ratio $= 1.002$, $p = 0.474$, across 224
bridge matrices at two model scales).
Rank dimensions in LoRA are rotationally symmetric: for any orthogonal
$R$, $B R^\top R A = BA$.
Topological labels on channels do not survive this symmetry without
explicit supervision.
The bridge learns \emph{that} mixing channels helps (FCC topology
produces $4.6\times$ more cross-channel coupling than cubic), but not
\emph{which} channels to prefer.

\paragraph{Contribution.}
We do not propose a better adapter.
We propose a new \emph{diagnostic}: a 36-parameter matrix that
summarizes the coupling structure training discovered between $A$ and
$B$, readable without inference or evaluation on held-out data.
The bridge provides a compact task fingerprint (1,008 parameters for
the optimal Q-projection subset---smaller than a single hidden state
vector in Qwen2.5-7B), a low-dimensional space for adapter composition
analysis, and a per-layer, per-module map of where task specialization
concentrates in the network.
To our knowledge, no prior work uses a structured intermediate
representation within an adapter as a diagnostic of adapter behavior.

The closest architectural precedent is LoRA-XS~\cite{banaei2024loraxs},
which places an $r \times r$ matrix between $A$ and $B$ but freezes
the outer projections to minimize parameter count.
Because $A$ and $B$ do not adapt to the task in LoRA-XS, the
intermediate matrix cannot serve as a diagnostic of what joint
$A$-$B$ learning discovered---it \emph{is} the entirety of learned
adaptation.
RhombiLoRA trains all three components jointly, enabling the bridge
to capture the emergent coupling pattern rather than carrying the full
adaptation burden.
Section~\ref{sec:related} discusses this distinction and other related
methods in detail.
